% Generated draft: requires mathematical and source review.
\documentclass[11pt]{article}
\usepackage[margin=25mm]{geometry}
\usepackage{fontspec}
\setmainfont{Latin Modern Roman}
\usepackage{amsmath,amssymb,mathtools,unicode-math}
\setmathfont{Latin Modern Math}
\usepackage{enumitem,needspace,xurl,hyperref}
\setlength{\parindent}{0pt}
\setlength{\parskip}{5pt}
\setlength{\emergencystretch}{3em}
\newcommand{\N}{\mathbb{N}}
\newcommand{\Z}{\mathbb{Z}}
\newcommand{\Q}{\mathbb{Q}}
\newcommand{\R}{\mathbb{R}}
\newcommand{\C}{\mathbb{C}}
\newcommand{\F}{\mathbb{F}}
\newcommand{\A}{\mathbb{A}}
\newcommand{\PP}{\mathbb{P}}
\newcommand{\E}{\mathbb{E}}
\newcommand{\eps}{\varepsilon}
\newcommand{\dd}{\,\mathrm{d}}
\begin{document}
\section*{UnsolvedMath ID: 2305018}
\textbf{Source problem code:} catalogue entry
\subsection{Definitions and mathematical statement}
Let $\mathbb D=\{z\in\C:|z|<1\}$, fix a real $0<\lambda<1$, and write
\[
f(z)=\lambda+\sum_{j\ge1}a_jz^j,\qquad
F_f(r)=\lambda+\sum_{j\ge1}|a_j|r^j\quad(0\le r<1).
\]
The source asks for the best radius $B(\lambda)$ such that $F_f(B(\lambda))\le1$ for every function in its intended class.

The literal extract assumes only that $f$ is analytic in $\mathbb D$. That is insufficient for a positive universal radius, since the linear coefficient is unrestricted. In particular the quoted positive values cannot be statements for the literal unbounded analytic class. A bound on the functions is missing.

Under the usual bounded-analytic interpretation, which must be confirmed rather than silently inserted, define
\[
\mathcal B_\lambda=\{f\text{ holomorphic on }\mathbb D:f(0)=\lambda,\ |f(z)|\le1\text{ on }\mathbb D\},
\]
\[
B(\lambda)=\sup\{r\in[0,1):F_f(r)\le1\text{ for every }f\in\mathcal B_\lambda\}.
\]
This makes ``best constant'' the largest permissible radius, with the inequality at the endpoint to be checked in the intended statement. With this explicit conditional interpretation, the requested unresolved range is $0<\lambda<1/2$. The source reports
\[
B(\lambda)=\frac{1}{1+2\lambda}\quad\left(\frac12\le\lambda<1\right),\qquad
B(0)=\frac1{\sqrt2},
\]
and
\[
B(\lambda)>\sqrt{\frac{1-\lambda}{2}}
\quad\left(0<\lambda<\frac12\right).
\]
The value at $\lambda=0$ is separately supplied endpoint information although the initial parameter range excludes zero. The strict lower bound is not a proposed exact value.

The supplied triage calls this Bombieri's $B(\lambda)$ problem and reports the full determination for $0<\lambda<1/2$ open as of the 2018 edition, with no later resolution located by that assessment. Its mention of de Branges and Bieberbach is contextual and supplies no reduction or theorem. No independent verification of the bounded-class correction, quoted formulas, or current status is claimed here.
\subsection{Short English statement}
Determine the largest radius at which the sum of absolute Taylor coefficients, weighted by that radius, stays at most one when the constant coefficient is fixed. The quoted positive radii require a missing bound on the analytic functions. Under the usual unit-bounded interpretation, the source leaves the interval $0<\lambda<1/2$ undetermined and supplies exact values elsewhere and a strict lower bound.
\subsection{Sources}
\begin{itemize}
\item Local source registry: \url{data/unsolvedmath/problems.json}, permanent UnsolvedMath ID $2305018$, source code catalogue entry, including the supplied OPEN-TRIAGE assessment. Dataset attribution: Ulam AI and UnsolvedMath Contributors; original-source attribution and rights are retained.
\item Hayman, \emph{Research Problems in Function Theory}, source-list date 2018, Problem 5.18: \url{https://arxiv.org/abs/1809.07200}. The metadata reports source-TeX extraction and July 29, 2026 access, with current-status and rights review flags.
\item The statement attributes the quoted radius facts to Bombieri, but its ``see Bombieri .'' reference is incomplete. The summary mentions de Branges and the Bieberbach conjecture without a precise relation or separate citation. No linked source was inspected for this draft.
\end{itemize}
\end{document}
